%! Unit Opener: Why Algebra Fluency Matters — Unit 1, Lesson 1.0 — Homework (Answer Key)
\documentclass[10pt]{article}
\usepackage{atda-article}
\usepackage{atda-key}   % pulls in -boxes; do NOT also load -boxes
\newcommand{\TallMath}[1]{$\displaystyle #1\rule[-1.4em]{0pt}{3.2em}$}
\setlength{\workrowsep}{8pt}

\begin{document}

\pageheader{Unit 1, Lesson 1.0}{Homework --- Answer Key}
% NAMESTRIP: no \namedateperiod here — mirrors the blank. Teacher-only prose goes
% in the lesson plan as \begin{teachernote}[Homework], never in a key.

\vspace{0.15in}

\boxguard[20]
\begin{notesbox}{Practice --- The Skills Unit 1 Is Built On}
\small
Items 1--3 are the diagnostic skills again with new numbers. Show your work.

\begin{enumerate}[leftmargin=1.6em, itemsep=6pt, topsep=4pt]

\item Simplify: $(2a^3b)(-5a^2b^4)$
\begin{work}
  (2a^3b)(-5a^2b^4) &= (2)(-5)\,a^{3+2}\,b^{1+4} \\
                    &= -10a^5b^5
\end{work}

\item Multiply and write in standard form: $(x-4)(x+9)$
\begin{work}
  (x-4)(x+9) &= x^2+9x-4x-36 \\
             &= x^2+5x-36
\end{work}

\item Factor out the greatest common factor: $18m^4-24m^3$
\begin{work}
  18m^4-24m^3 &= 6m^3(3m)-6m^3(4) \\
              &= 6m^3(3m-4)
\end{work}

\end{enumerate}
\end{notesbox}

\vspace{0.10in}

\boxguard[24]
\begin{scenariobox}[In Context --- Choosing the Right Form]{royal}
\small
A softball is tossed straight up from a raised mound. Its height above the ground, in feet,
$t$ seconds after release is
\[ h(t) = -16t^2+32t+48 \qquad\text{and also}\qquad h(t) = -16(t-3)(t+1). \]

\begin{enumerate}[leftmargin=1.6em, itemsep=6pt, topsep=4pt, start=4]

\item Which form tells you the height the ball was released from, and what is that height?

\ansline{The expanded form --- its constant term gives $h(0)=48$ feet.}\\

\item Which form tells you when the ball hits the ground, and when is that? (Remember: one of the
two zeros does not describe the toss.)

\ansline{The factored form --- the zeros are $t=3$ and $t=-1$, so it lands at 3 seconds.}\\

\item Verify that the two forms really are equivalent by expanding the factored form.
\begin{work}
  -16(t-3)(t+1) &= -16\big(t^2+t-3t-3\big) \\
                &= -16\big(t^2-2t-3\big) \\
                &= -16t^2+32t+48
\end{work}

\end{enumerate}
\end{scenariobox}

\vspace{0.10in}

\boxguard[14]
\begin{scenariobox}[A First Look Ahead]{goldacc}
\small
\begin{enumerate}[leftmargin=1.6em, itemsep=6pt, topsep=4pt, start=7]

\item You deposit \$500 in an account that grows by $4\%$ each year. After $t$ years the balance
in dollars is $A(t)=500(1.04)^t$. Find $A(3)$, rounded to the nearest cent, then write one
sentence saying what it means.
\begin{work}
  A(3) &= 500(1.04)^3 \\
       &= 500(1.124864) \\
       &\approx 562.43
\end{work}
\ansline{After 3 years the account holds about \$562.43, up \$62.43 from the deposit.}

\end{enumerate}
\end{scenariobox}

\vspace{0.10in}

\boxguard[16]
\begin{extensionbox}
\small
\textbf{The number trick.} Pick any number. Add $7$. Double the result. Subtract $4$. Divide by
$2$. Now subtract the number you started with. You get $5$ --- every time.

Let $n$ be the number you picked, write an expression for each step, and simplify to explain why
the answer is always $5$.
\begin{work}
  \dfrac{2(n+7)-4}{2}-n &= \dfrac{2n+10}{2}-n \\
                        &= (n+5)-n \\
                        &= 5
\end{work}
\end{extensionbox}

\vspace{0.10in}

\begin{remindbox}
\small
\textbf{Next: Lesson 1.1 --- Simplifying Algebraic Expressions and Exponent Rules}
(\texttt{A2.EO.3a}). Every ``expand and compare'' you did today runs on the exponent rules and
the distributive property. Lesson 1.1 makes them fast and automatic --- and those same rules
come back in Unit 3 as the properties of logarithms.
\end{remindbox}

\end{document}
